\title{Bewertungsraster Factsheet\\ \large Informatik \& Gesellschaft}

\maketitle

\section*{Hinweise zum Factsheet}

\textbf{Bewertungskriterien}:
\begin{itemize}
    \item (2 Pt) Inhaltlich korrekte und umfassende Zusammenfassung des Kapitels
    \item (1 Pt) Mindestens eine selbst erstellte \textbf{Illustration} (Bild, Diagramm, Grafik, usw.), die Ihr Factsheet sinnvoll ergänzt.
    \item (1 Pt) Formale Kriterien
    \begin{itemize}
        \item Abgabeformat: PDF
        \item Dateiname: ``\texttt{[Klassenname] - K[Kapitelnummer] [Kapitelname] Factsheet.pdf}'' (z.B. ``\texttt{2d - K2 Daten Factsheet.pdf}'')
        \item 3-4 A4-Seiten
        \item Korrekte Sprache, vollständige Sätze
        \item Schriftgrösse: 11 Punkte
        \item Vor- und Nachnamen aller Gruppenmitglieder im PDF enthalten
    \end{itemize}
    \item (1 Pt) Formulieren Sie am Ende Ihrer Zusammenfassung zwei interessante, reichhaltige Diskussions-Fragen (die nicht einfach mit ja / nein beantwortet werden können) oder provokativ-pointierte Aussagen, die zur Diskussion anregen. Versuchen Sie, Fragen zu formulieren, die Bezug zu Ihrer eigenen Lebenswelt (Interessen, Hobbies etc.) haben.
\end{itemize}


Das Fact Sheet enthält zudem spannende Fragestellungen, welche die Grundlage für die Gruppendiskussionen bilden

Geben Sie im Anhang (zählt nicht zu den 3-4 Seiten) an, wer für welche(n) Unterabschnitt(e) zuständig war.

\vspace{1em}

Das Factsheet ist spätestens eine Woche vor dem Vortrag abzugeben.


\newpage

\section*{Bewertung Factsheet: Justiz und Algorithmen}

\textbf{Autorinnen: Ladina, Maha, Minna, Salome, Aysima}

\begin{itemize}
    \item \textbf{Inhaltliche Zusammenfassung (2/2 Punkte):} Das Factsheet bietet eine sehr umfassende und verständliche Zusammenfassung aller relevanten Aspekte des Themas. Die wichtigsten Probleme, Chancen und Risiken werden differenziert dargestellt und mit Beispielen belegt. Die Unterkapitel orientieren sich sinnvoll an den Leitfragen aus der Aufgabenstellung.
    \item \textbf{Illustrationen (1/1 Punkt):} Es sind sinnvolle und anschauliche Visualisierungen enthalten, die das Verständnis für neuronale Netze, Entscheidungsbäume und Zusammenhänge im Thema fördern. Die Karte ist informativ und passend gewählt.
    \item \textbf{Formale Kriterien (1/1 Punkt):} Das Dokument erfüllt alle formalen Vorgaben (korrekte Sprache, vollständige Sätze, richtige Länge und Formatierung, alle Namen enthalten).
    \item \textbf{Diskussionsfragen (1/1 Punkt):} Am Ende werden zwei interessante, offene Diskussionsfragen gestellt, die zur kritischen Auseinandersetzung mit dem Thema anregen und einen Bezug zur Lebenswelt ermöglichen.
\end{itemize}

\textbf{Gesamtpunktzahl: 5/5 Punkten}
